\documentclass[11pt]{article}
\usepackage[margin=1in]{geometry}
\usepackage{amsmath}
\usepackage{amssymb}
\usepackage{graphicx}
\usepackage{multicol}

\setlength{\parindent}{0pt}
\setlength{\parskip}{1\baselineskip}

\title{UnlearnHead: Inference-Time Continual Unlearning with Lightweight Logit Heads}

% \author{Juan Belieni}

\begin{document}

\maketitle

% \section{Preliminaries}

% \subsection{Autoregressive Language Modeling}

% Let $F_\theta$ be a transformer language model with hidden dimension $d$ and vocabulary $\mathcal{V}$. Given a token prefix $x_{\leq t}$, the transformer produces a contextual representation $h_t=F_\theta(x_{\leq t})\in\mathbb{R}^d$. A language-model head $W_0\in\mathbb{R}^{|\mathcal{V}|\times d}$ maps this representation to logits $z_t^0=W_0h_t$, which define the next-token distribution
% \begin{equation}
%     p_0(v\mid x_{\leq t})
%     =\frac{\exp(z_t^0(v))}
%     {\sum_{u\in\mathcal{V}}\exp(z_t^0(u))},
%     \qquad v\in\mathcal{V}.
% \end{equation}
% Generation proceeds autoregressively by selecting or sampling a token from this distribution and appending it to the prefix. This separation between representation computation and next-token scoring is central to our setting: we retain the representation produced by the pretrained model and intervene only in the logits used for decoding.

% \subsection{LLM Unlearning}

% LLM unlearning seeks to reduce a model's ability to reproduce designated information while preserving its behavior on information that should remain available. We denote the examples associated with removal request $i$ by the forget set $D_i^F$ and the examples used to preserve general utility by the fixed retain set $D^R$. Existing approaches commonly modify model parameters by optimizing against the forget data, for example through gradient-based knowledge unlearning \cite{jang2022knowledge} or negative-preference objectives \cite{zhang2024npo,fan2024simnpo}. This creates a balance between forgetting efficacy, retained utility, robustness, and computational cost \cite{cha2024robust,yuan2024closer}.

% This balance is particularly important in the continual setting, where requests arrive as a sequence $D_1^F,\ldots,D_K^F$. Repeatedly updating a shared parameter set can make the cost grow with the number of requests and can allow later updates to interfere with earlier behavior. We therefore consider a modular alternative: the base language model remains unchanged, while each request contributes a small component that can be added independently.

% Unlearning also cannot be assessed solely by a loss decrease or by failure on one set of prompts. Benchmarks such as TOFU \cite{maini2024tofu} and MUSE \cite{shi2024muse} evaluate forgetting together with model utility and related behavioral criteria, while recent work has highlighted weaknesses in current benchmarks and the possibility that apparently forgotten information remains recoverable under alternative queries or decoding strategies \cite{thaker2024weak,doshi2024blackbox,reisizadeh2025leakk}. Accordingly, our use of \emph{unlearning} refers to an inference-time behavioral intervention; it does not claim that the underlying information has been removed from the frozen model parameters.

% \subsection{Contrastive Decoding and Lightweight Adaptation}

% Decoding-time control provides a way to alter generation without changing the base model. Prior methods steer a language model by combining its predictions with scores from discriminators, experts, or anti-experts \cite{krause2020gedi,yang2021fudge,liu2021dexperts}. More generally, contrastive decoding changes token preferences by comparing predictive distributions and using their difference as a decoding signal \cite{li2022contrastive}. UCD applies this principle to LLM unlearning: an original model is contrasted with auxiliary models trained on forget and retain data, and the resulting signal is used to suppress forget-associated continuations at inference time \cite{suriyakumar2025ucd}.

% The effectiveness of this strategy suggests that unlearning need not always be expressed as a permanent update to the base model. Its main cost, however, depends on how the auxiliary distributions are produced. Full auxiliary language models require additional parameter storage and transformer computation. Parameter-efficient methods show that useful task-specific changes can instead be represented by small modules \cite{houlsby2019adapters}, including low-rank updates \cite{hu2021lora}. We combine these two ideas: contrast supplies the unlearning signal, while low-rank logit heads produce that signal directly from a single shared hidden state.


\section{Method}

\subsection{Unlearning Without Modifying the Language Model}

Continual unlearning poses a practical tension. Updating the full language model after every removal request is costly, and repeated updates may progressively damage capabilities unrelated to the requested information. We therefore leave the pretrained transformer and its language-model head unchanged. Instead of attempting to erase information from the model parameters, we learn small request-specific modules that identify how the model's predictions should be corrected when forgotten information becomes relevant.

For a prefix $x_{\leq t}$, the frozen transformer is evaluated once to produce its final hidden representation, after which the original language-model head produces the base logits:
\begin{equation}
    h_t = F_\theta(x_{\leq t}),
    \qquad
    z_t^0 = W_0 h_t.
\end{equation}
All subsequent computations operate on the shared representation $h_t$. Thus, adding an unlearning request does not introduce another transformer forward pass; it only adds a small computation on top of a representation that the model already produces during generation.

\subsection{Learning What Distinguishes Forget Data}

The method must distinguish tokens that are characteristic of the information to be forgotten from tokens that are generally plausible in the same context. To obtain this distinction, we compare two auxiliary views of the frozen representation. A request-specific forget head $H_i^F$ is trained on $D_i^F$, while a shared retain head $H^R$ is trained once on the fixed retain set $D^R$. Each forget head is paired with this same retain head to isolate the preferences associated with its requested information from those supported by data that should remain available.

Both heads are low-rank residual maps rather than complete vocabulary projections:
\begin{equation}
    \delta_i^F(h)=B_i^F A_i^F h,
    \qquad
    \delta^R(h)=B^R A^R h,
\end{equation}
where $A\in\mathbb{R}^{r\times d}$ projects a hidden state of dimension $d$ to a small rank $r$, and $B\in\mathbb{R}^{|\mathcal{V}|\times r}$ maps it to vocabulary space. These residuals define the specialized logits
\begin{equation}
    z_i^F=z^0+\delta_i^F(h),
    \qquad
    z^R=z^0+\delta^R(h).
\end{equation}
Because both heads reuse the base logits, they need only learn the adjustments required to specialize the frozen model toward their respective datasets.

The heads are trained by ordinary next-token prediction:
\begin{equation}
    \mathcal{L}_F^{(i)}=-\sum_t \log p_i^F(y_t\mid y_{<t}),
    \qquad
    \mathcal{L}_R=-\sum_t \log p^R(y_t\mid y_{<t}).
\end{equation}
The forget head consequently assigns greater probability to continuations associated with request $i$, whereas the retain head captures preferences supported by the fixed reference data. This direct training avoids both a teacher model and target-logit distillation.

The useful quantity is not either distribution in isolation, but their contrast. We define the logit contrast for request $i$ as
\begin{equation}
    \Delta_{i,t}
    = z_i^F-z^R
    =\delta_i^F(h_t)-\delta^R(h_t),
\end{equation}
where the base logits cancel because they are shared by the two specialized heads. For any vocabulary tokens $u$ and $v$, a larger value of $\Delta_{i,t}(v)$ relative to $\Delta_{i,t}(u)$ means that the forget-specialized distribution favors $v$ over $u$ more strongly than the retain distribution does. Moreover, the log-probability ratio differs from this logit contrast only by a scalar that is constant across the vocabulary:
\begin{equation}
    \log p_i^F(v\mid x_{\leq t})
    -\log p^R(v\mid x_{\leq t})
    =\Delta_{i,t}(v)+c_{i,t}.
\end{equation}
Because softmax is invariant to adding the same scalar to every logit, subtracting the normalized log-probability contrast or subtracting $\Delta_{i,t}$ produces the same corrected next-token distribution. We therefore use the directly computed logit contrast during decoding. Appendix~\ref{app:logit-contrast} gives the derivation. The normalized probabilities are still used when computing the routing score below, since their normalization terms vary across prefixes and cannot in general be omitted from a sequence likelihood ratio.

\subsection{From Contrast to an Unlearned Distribution}

Having identified the directions in logit space that favor the forget data, we reverse those directions during decoding. The correction associated with request $i$ is
\begin{equation}
    z_{i,t}^U=z_t^0-\alpha_i(q)
    \left(\delta_i^F(h_t)-\delta^R(h_t)\right),
\end{equation}
where the prompt-dependent weight $\alpha_i(q)\geq0$ is defined by the routing mechanism below. Irrelevant requests receive zero weight, while requests judged relevant share a fixed global correction budget.

The corrected logits are passed to a standard softmax and can be decoded using greedy, top-$k$, nucleus, or any other conventional sampling procedure.

\subsection{Calibrated Likelihood-Ratio Routing}

The contrast between the auxiliary heads provides a natural measure of when a correction should be applied. The request-specific forget head represents the distribution associated with request $i$, while the shared retain head provides a common reference distribution. Together they act as a generative relevance detector for the request. For a prompt $q=(x_1,\ldots,x_T)$, we compute the average log-likelihood ratio
\begin{equation}
    s_i(q)=\frac{1}{T}\sum_{t=1}^{T}
    \left[
    \log p_i^F(x_t\mid x_{<t})
    -\log p^R(x_t\mid x_{<t})
    \right].
\end{equation}
A large score indicates that the prompt is more characteristic of the forget distribution than of the fixed retain distribution. The score has a direct probabilistic interpretation and is evaluated for every stored request using the same frozen transformer hidden states. The shared retain head is evaluated once, and each request adds only the computation of its lightweight forget head.

Raw scores are not directly comparable across requests because independently trained forget heads may induce different score distributions. We calibrate each request on a fixed, non-sensitive set of unrelated prompts $D_{\mathrm{cal}}$. Its empirical negative-score distribution is
\begin{equation}
    F_i^-(s)
    =\frac{1}{|D_{\mathrm{cal}}|}
    \sum_{q\in D_{\mathrm{cal}}}
    \mathbb{I}\left[s_i(q)\leq s\right].
\end{equation}
Given a target false-positive rate $\varepsilon$, the request-specific activation threshold is the empirical $(1-\varepsilon)$-quantile
\begin{equation}
    \tau_i
    =Q_{1-\varepsilon}
    \left(\left\{s_i(q):q\in D_{\mathrm{cal}}\right\}\right).
\end{equation}
This calibration directly associates the threshold with the desired routing behavior: under the calibration distribution, request $i$ activates on at most approximately an $\varepsilon$ fraction of unrelated prompts. It also accommodates request-specific and non-Gaussian score distributions without imposing a common score scale.

The router then applies the binary gate
\begin{equation}
    \gamma_i(q)=\mathbb{I}
    \left[s_i(q)>\tau_i\right].
\end{equation}
Multiple gates may activate simultaneously when a prompt concerns more than one removal request. Let
\begin{equation}
    N(q)=\sum_{j=1}^{K}\gamma_j(q)
\end{equation}
be the number of active requests. We assign equal weight to all active requests:
\begin{equation}
    \alpha_i(q)
    =\lambda\frac{\gamma_i(q)}{\max\{1,N(q)\}},
    \qquad 0<\lambda\leq1,
\end{equation}
where $\lambda$ is a global correction-strength parameter. If no gate activates, then every $\alpha_i(q)$ is zero and the base model is unchanged. Otherwise, the active weights sum to $\lambda$; choosing $\lambda<1$ reserves part of the correction budget for the unchanged base behavior. Equal normalization also prevents the total intervention strength from increasing merely because several requests activate. The target false-positive rate $\varepsilon$ controls the trade-off between unnecessary activation and failure to detect relevant prompts.

\subsection{Continual Operation}

Each new removal request extends the system rather than rewriting it. The shared retain head $H^R$ is trained once on $D^R$ and remains fixed. Given a new forget set $D_i^F$, we freeze the base model and $H^R$, train $H_i^F$, evaluate $s_i$ on $D_{\mathrm{cal}}$, and compute the scalar threshold $\tau_i$. The resulting request state is
\begin{equation}
    R_i=\left(H_i^F,\tau_i\right),
\end{equation}
which is added to the head bank. After $K$ requests, the full auxiliary state consists of the shared retain head and the request states $\{R_i\}_{i=1}^K$. Because each request contributes an independent forget head and threshold, adding it leaves all existing components fixed. At inference time, the calibrated gates determine which requests are relevant and their corrections are combined as
\begin{equation}
    z_t^U=z_t^0-\sum_{i=1}^{K}
    \alpha_i(q)
    \left(\delta_i^F(h_t)-\delta^R(h_t)\right).
\end{equation}
Because inactive requests have zero weight, generation evaluates contrastive corrections only for the active heads. This construction preserves the central computational benefit of the method: generation requires one transformer evaluation, one ordinary language-model-head evaluation, and only the active low-rank projections.

This organization makes the same learned contrast serve two roles: it determines the direction of the logit correction and measures whether that correction is relevant to the prompt. Once $\tau_i$ has been computed, it summarizes the calibration prompts for routing purposes and is stored with the corresponding request state.

The routing mechanism can be evaluated by measuring its false-positive rate on unrelated prompts and its recall on prompts related to each forget request. Robustness requires extending this evaluation to paraphrases and indirect queries, while continual operation additionally calls for measuring calibration as the number of requests grows and testing prompts for which several gates should activate simultaneously.

The design rests on one central assumption. Although the output behavior is altered at inference time, the representation itself remains fixed. The hidden states must therefore already expose enough information for lightweight heads to separate forget-specific preferences from general language-model behavior. This makes the approach less expressive than methods that independently fine-tune full auxiliary models, but substantially cheaper to extend and execute. Whether this restricted representation is sufficient to achieve reliable forgetting while retaining unrelated capabilities is the main empirical question evaluated by the method.

\bibliographystyle{plain}
\bibliography{refs}

\appendix

\section{Equivalence of Log-Probability and Logit Contrasts}
\label{app:logit-contrast}

Fix a request $i$, decoding step $t$, and hidden state $h_t$. Let
\begin{equation}
    Z_{i,t}^F=\sum_{u\in\mathcal{V}}\exp\!\left(z_i^F(u)\right),
    \qquad
    Z_t^R=\sum_{u\in\mathcal{V}}\exp\!\left(z^R(u)\right)
\end{equation}
be the partition functions of the forget and retain heads. For any token $v\in\mathcal{V}$,
\begin{align}
    \log p_i^F(v\mid x_{\leq t})
    -\log p^R(v\mid x_{\leq t})
    &=z_i^F(v)-\log Z_{i,t}^F-z^R(v)+\log Z_t^R \\
    &=\Delta_{i,t}(v)+c_{i,t},
\end{align}
where
\begin{equation}
    c_{i,t}=\log Z_t^R-\log Z_{i,t}^F
\end{equation}
does not depend on $v$. The two specialized heads share the same base logits, so
\begin{equation}
    \Delta_{i,t}
    =z_i^F-z^R
    =\left(z_t^0+\delta_i^F(h_t)\right)
     -\left(z_t^0+\delta^R(h_t)\right)
    =\delta_i^F(h_t)-\delta^R(h_t).
\end{equation}

To see why the constant can be omitted during decoding, let $a=\alpha_i(q)$ and define the corrected logits using the exact log-probability contrast. They are
\begin{align}
    \widetilde z_{i,t}^U(v)
    &=z_t^0(v)-a
    \left[\Delta_{i,t}(v)+c_{i,t}\right] \\
    &=\left[z_t^0(v)-a\Delta_{i,t}(v)\right]
      -a c_{i,t}.
\end{align}
The last term is identical for every vocabulary token. Hence
\begin{equation}
    \operatorname{softmax}\!\left(\widetilde z_{i,t}^U\right)
    =\operatorname{softmax}\!\left(z_t^0-a\Delta_{i,t}\right).
\end{equation}
Thus, for the linear contrastive correction used in this work, the difference between the residual logits is an exact representative of the log-probability contrast for purposes of the resulting next-token distribution. This equivalence relies on applying the contrast linearly before softmax; it does not generally hold after elementwise nonlinear transformations.

\end{document}
